\chapter{Empirical Convergence: A Map of the State of the Art}

\epigraph{\emph{The evidence does not yet form a theory, but it does form a pattern.}}{}

The previous part of this manuscript defined the conceptual object of HIT and stated the hierarchy of hypotheses that organize the framework. The present chapter changes register. Its task is not to introduce new ontological claims, nor to defend the experimental program in detail, but to map the state of knowledge that already surrounds the problem from several directions at once. The basic contention is simple. A substantial portion of what HIT seeks to name has already been discovered in fragments. The difficulty is not that the literatures are empty. The difficulty is that they are distributed across different disciplines, scales, methods, and vocabularies. What one field describes as resonance, another describes as periodicity, coupling, roughness, entrainment, harmonicity, recurrence, niche partitioning, code, or prediction. The objects are not identical, and the mechanisms are not always the same, but the family resemblance is too strong to ignore.

This chapter therefore reads the state of the art as a problem of convergence rather than as a single settled doctrine. It asks four questions of each domain it surveys: what has this field discovered, what vocabulary does it use to describe those discoveries, what remains outside its normal frame, and how does its local insight bear on the broader problem of harmonic information? That is the level at which the literature becomes illuminating for HIT. The chapter is not exhaustive in the archival sense; no single review of reasonable length could be. It is selective in a more demanding way. It assembles the strands most relevant to the central claims of the framework while preserving their disciplinary specificity. That preservation matters. The chapter does not collapse the fields into one another, and it does not mistake analogy for identity. Its purpose is to show that there is already enough empirical and conceptual material on the table to justify a coordinated research program.

\section{Neuroscience: resonance, coupling, and neural organization}

The neuroscientific literature provides one of the clearest indications that structured relations among frequencies are more than descriptive conveniences, which is why this field matters so much for HIT. Within one of the best-instrumented empirical domains available, the issue first becomes visible through work on neural resonance and tonal processing. Large and colleagues argued that tonal organization is not best understood as passive response to isolated frequencies, but as the dynamic entrainment of oscillatory systems to patterned temporal relations \parencite{large_2005, large_2012}. That shift matters because it displaces a stimulus-centered account in favor of a relation-centered one. Tonality, in this view, is not merely a label applied after the fact to a collection of notes. It emerges from how neural systems lock onto periodic structure across time. Even when the specific models remain debated, the broader lesson is stable: the nervous system appears to treat some temporal and harmonic organizations as privileged constraints rather than neutral inputs.

This line of work is reinforced by studies of subcortical and cortical sensitivity to consonance and harmonic structure. Bidelman and collaborators showed that consonant relations can be reflected in early auditory processing and in hierarchical neural organization, suggesting that the discrimination of harmonic order is not exhausted by late cultural interpretation alone \parencite{bidelman_2009, bidelman_2011}. The force of these findings lies less in any final claim about universality than in the fact that the nervous system appears able to register patterned intervallic organization before one reaches the level of explicit musical judgment. Other neural studies, including intracranial and EEG-based work on tonal discrimination and consonance processing, move in the same direction: harmonic organization is repeatedly treated as a measurable feature of neural response rather than as a purely symbolic overlay.

Cross-frequency coupling broadens this picture beyond music in the narrow sense. Canolty and Knight showed that relations among oscillatory bands can operate as mechanisms of large-scale integration in the brain rather than as incidental correlations \parencite{canolty_2010}. Here the vocabulary changes, but the structural problem is familiar. What matters is not only the existence of oscillation, but the patterned coordination among oscillations at different scales. This is one of the reasons the neuroscientific literature remains so important for HIT. It shows that the question of harmonic information does not begin and end with heard intervals. It extends into the way nervous systems organize communication among their own rhythmic processes. Ratios, couplings, and phase relations become relevant because they modulate integration, not because they satisfy an aesthetic preference.

The work associated with Varela extends this point from local coupling to large-scale neural integration. In the Brainweb synthesis and related studies, distributed neural ensembles are described as transiently coordinated through phase synchronization rather than by means of a single central controller \parencite{varela_2001_brainweb, rodriguez_1999}. That picture matters for HIT because it treats resonance as a mechanism of integration and not as a decorative metaphor. Coherence emerges when distributed populations converge on temporally compatible relations, often with strikingly musical language but explicitly without a conductor. Singer's work on neuronal synchrony reinforces the same lesson from a nearby angle: relations of timing can define what belongs together and what does not within ongoing neural organization \parencite{singer_1999}. None of these studies proves HIT. They do show that the nervous system already uses relation, phase, and resonant coupling as explanatory primitives in one of the best-instrumented empirical domains available.

Recent work on traveling waves sharpens the point further. Jacobs and colleagues argued that eigenfrequency relations and wave geometry help structure cortical organization in ways that are neither purely local nor purely metaphorical \parencite{jacobs_2025}. The significance of this work for HIT lies not in proving that the brain \enquote{is musical,} but in demonstrating that relations among frequencies can bear on geometry, communication, and organization at once. Neural systems do not merely host oscillations; they appear to exploit patterned relations among them. What neuroscience often leaves less developed is the broader comparative horizon. The field excels at showing how frequency relations participate in neural integration, but it usually stops short of connecting those findings to analogous structures in ecology, dynamical systems, or informational theory. It sees one essential part of the picture with unusual clarity, but it does not yet provide the whole map.

\section{Psychoacoustics and tonal perception}

If neuroscience shows that frequency relations matter for neural organization, psychoacoustics shows how richly the sonic domain has already elaborated the problem at the perceptual level; without it, the map would lose its most precise historical vocabulary, even if the field often treats the auditory case as the natural border of the problem. The classic starting point remains Helmholtz, whose work on tonal sensation displaced purely arithmetic accounts of consonance by grounding the phenomenon in the interaction of overtones and the physiology of hearing \parencite{helmholtz_1863}. Plomp and Levelt later gave this intuition one of its canonical formulations through critical-band theory, showing that roughness and interference among partials are central to why some intervals are experienced as more consonant than others \parencite{plomp_1965}. These traditions remain decisive because they refuse both naive numerology and purely cultural arbitrariness. They insist that consonance is a structured phenomenon arising from the interaction of ratios, spectra, and auditory mechanisms.

Later work complicated the picture in productive ways. Roederer and Oxenham both emphasized that pitch perception, harmonicity, and tonal organization cannot be reduced to one simple cue or one simple model \parencite{roederer_2008, oxenham_2012}. The auditory system does not infer organization from a single fixed principle; it integrates temporal, spectral, and contextual information under changing conditions. This complication does not weaken the field's importance for HIT. On the contrary, it helps clarify why harmonic information should be treated relationally rather than as a property of bare number. The psychoacoustic archive shows, again and again, that structured proportional relations matter, but always in concert with the modal constitution of the signal and the constraints of the receiver.

Another important branch of this literature connects tonal preference to the statistics of vocal production and the biology of communication. Schwartz and colleagues argued that some intervallic regularities in music reflect the statistical structure of voiced speech, while Bowling and Purves developed related arguments linking consonance to biologically familiar harmonic organization \parencite{schwartz_2003, bowling_2015}. These lines of work remain contested, and they should remain contested. The relevant point here is not that they settle the origin of consonance, but that they frame harmonic organization as something more than a local cultural accident. In this respect the psychoacoustic literature becomes one of the main historical laboratories for HIT. It gives the problem an unusually rich vocabulary, even where its dominant framing remains largely auditory.

That auditory concentration is both the field's strength and its limit. Psychoacoustics often sees the phenomenon with remarkable precision when the signal is sonic, the task is perceptual, and the subject is human or near-human. But the same precision can narrow the horizon. The literature tends to localize the question within hearing rather than ask whether the structures it studies might belong to a larger family of resonant or informational constraints. HIT does not erase the field's local specificity; it depends on it. At the same time, the project must read psychoacoustics as revealing a particularly legible case of a broader problem rather than as defining the borders of that problem once and for all.

\section{Acoustic ecology and biosemiotics}

The ecological and biosemiotic literatures shift attention from isolated perceptual judgments to organized environments and living systems, and they matter because HIT would remain too narrow if harmonic organization were framed only as laboratory perception rather than also as patterned coexistence, signaling, and interpretation in living environments. Krause's work on biophony and the acoustic niche hypothesis is central here because it treats soundscapes as structured ecologies rather than undifferentiated acoustic mass \parencite{krause_2013}. In such environments, organisms do not simply emit signals into a neutral background. They occupy, avoid, and negotiate spectral and temporal space in patterned ways. This is important for HIT because it introduces a form of order that is neither purely musical nor reducible to laboratory psychophysics. The key fact is patterned coexistence. Some bands, intervals, and temporal windows become functionally significant because they support coexistence, communication, or avoidance under ecological constraints.

Buxton and colleagues extended part of this ecological insight into the domain of human well-being, showing that exposure to natural soundscapes can have measurable benefits for health and restoration \parencite{buxton_2021}. Snowdon's work on animal signaling and welfare broadens the point further by treating patterned acoustic environments as biologically consequential rather than ornamentally pleasant \parencite{snowdon_2021}. These traditions do not yet formulate a theory of harmonic information, but they do reveal something decisive for its plausibility: living systems are often embedded in non-random signal ecologies where patterned organization matters for function, not only for representation.

Autopoiesis enters fruitfully at this point because it supplies a vocabulary for the living unity that produces, interprets, and is transformed by such patterned fields. Hoffmeyer and Barbieri describe signs, codes, and semiotic organization with great force, but Maturana and Varela describe the kind of system for which any such code can matter in the first place: a unity whose continued organization depends on the maintenance of a relational network of processes \parencite{maturana_varela_1980}. Read through that lens, ecological patterning is not only a matter of better signaling. It is part of the ongoing conditions under which a living system preserves congruence with its environment. Di Paolo's notion of precarious autonomy makes the normative edge of the problem explicit. A living system is not merely organized; it is vulnerable, and the loss of viable organization matters because the system can disintegrate \parencite{dipaolo_2005}. This is one reason harmony becomes theoretically important here. Its value need not be aesthetic before it is organizational.

Biosemiotics articulates this problem from another direction. Hoffmeyer and Barbieri both insist, in different ways, that living systems cannot be adequately described as mere exchanges of matter and energy; they are also systems of sign, code, and interpretive organization \parencite{hoffmeyer_2008, barbieri_2008}. This does not mean that every harmonic pattern is automatically semiotic. It means that when recurring relational forms become functionally efficacious for living systems, the language of sign and interpretation becomes relevant. The affinity with HIT is clear. A patterned ratio or recurrence structure matters only if it can make a difference for some system. That insistence brings the ecological and biosemiotic literatures into a productive relation with the more formal accounts of information discussed later in the chapter.

Ethnographic work such as Feld's study of Kaluli acoustemology adds a final dimension. In such work, forest sound, social life, orientation, and musical form are not separable layers but mutually informing structures \parencite{feld_1982}. This does not provide a universal theory, nor should it be made to do so. What it shows is that patterned environments can be lived as organized worlds rather than merely measured as spectra. Acoustic ecology and biosemiotics therefore illuminate aspects of the problem that both psychoacoustics and nonlinear dynamics often underemphasize: that patterned organization can be ecological, semiotic, and embodied at once. What these literatures often lack is a sharper formal vocabulary for recurrence, ratio structure, and dynamic stability. They see living organization vividly, but they do not always possess the mathematical language needed to connect that organization to the rest of the convergent map.

\section{Dynamical systems, recurrence, and synchronization}

If acoustic ecology gives the problem an environmental and living form, nonlinear dynamics enters the map at a specific juncture because it provides the formal skeleton that many neighboring literatures presuppose without fully stating. Kuramoto's model of coupled oscillators remains foundational because it shows, with unusual economy, that some relations among oscillatory systems are dynamically privileged \parencite{kuramoto_1984}. Phase locking does not emerge under arbitrary conditions with equal ease; simple rational relations often define wider and more stable coordination regimes. The later literature on synchronization, from Strogatz to Pikovsky and Acebron, made this point central to the understanding of complex systems more broadly \parencite{strogatz_2000, pikovsky_2001, acebron_2005}. In this literature, ratio simplicity is not treated as a mystical property. It is treated as a recurring structural condition under which coordination becomes more stable, more robust, or more widespread.

The importance of recurrence analysis lies in the way it translates these intuitions into observable structure. Eckmann and colleagues introduced recurrence plots as a way of visualizing returns in dynamical systems, and Marwan later consolidated recurrence quantification analysis into a substantial analytical framework \parencite{eckmann_1987, marwan_2007}. For HIT, Trulla and collaborators are especially important because they connect recurrence structure directly to consonance-like organization, showing that simpler harmonic relations produce more stable and legible recurrence patterns in the time domain \parencite{trulla_2018}. This matters not because it solves the whole problem of consonance, but because it gives one strong answer to a narrower question: what kind of measurable dynamical profile tends to accompany simple harmonic relations? The answer is not aesthetic approval, but patterned return.

The literature on biological oscillators strengthens this bridge between formal dynamics and living organization. Glass and Mackey showed that physiological systems often exhibit ratio-structured coordination and mode-locking behavior that is neither accidental nor confined to the laboratory \parencite{glass_1988}. Schroeder's work on devil's staircases and Cartwright's work on pitch as mode-locking point in a similar direction from different angles: stable step-like structures, locking regimes, and perceptual organization can emerge from common dynamical principles \parencite{schroeder_1991, cartwright_2001}. What this family of work contributes is a vocabulary for discussing harmonic privilege without immediately invoking cultural tradition, auditory preference, or metaphysical language. It shows that the problem can be formalized as one of recurrent stability, locking, and attractor structure.

Even geometric visualizations such as Lissajous figures belong to this formal family. Gallozzi and Strollo's work makes clear that simple ratios generate topologically cleaner and more easily closed trajectories than more complex relations do \parencite{gallozzi_2023}. Once again, the point is not that visual simplicity settles the ontology of consonance. It is that different domains keep discovering that simple proportional relations produce a recognizable economy of recurrence and organization. This is one of the places where the trans-physical framing of HIT becomes especially persuasive. Nonlinear dynamics does not need to speak the language of music for its findings to matter to the same problem.

The same point becomes even sharper when one shifts from curve shape alone to nodal and singular structures in interference fields. From the older archive of Chladni plates and cymatic visualization to Berry and Dennis's work on phase singularities in monochromatic waves, one finds repeated evidence that interference does not merely generate pretty images but physically constrained topologies with stable structural features \parencite{berry_2001}. This matters for HIT because it widens the empirical meaning of harmonic organization. A simple ratio can manifest not only as a sounding interval or a recurrent trajectory, but also as a visible field organization with its own nodes, closures, dislocations, and transformations. The geometry does not replace the dynamics. It makes part of the dynamics legible.

Controlled interference media sharpen this point further because they allow one to observe not only pattern formation but transitions between pattern regimes. When harmonic, rational, and incommensurable relations are swept under stable boundary conditions, what becomes visible is not a binary opposition between order and disorder, but a family of differences in closure, drift, and topological stability. This is one of the reasons such systems matter to HIT. They make it possible to ask, in a disciplined way, which aspects of a pattern belong to the excitatory input, which belong to the resonant structure of the medium, and which emerge only from their interaction. In that sense, controlled interference belongs to the same empirical map as recurrence analysis: it renders relational organization measurable rather than merely suggestive.

The limit of this literature is equally clear. It often sees the formal organization with extraordinary sharpness while remaining comparatively silent about meaning, embodiment, ecology, or perception. It can show that some relations are dynamically privileged without yet telling us why organisms, nervous systems, or cultures care about them in the ways they do. That is not a flaw in the field; it is simply the boundary of its normal questions. For HIT, however, the lesson is decisive. The dynamical systems literature provides the strongest formal argument that simple harmonic relations are not arbitrary conveniences, but real candidates for explanatory priority.

\section{Comparative, cross-cultural, and historical evidence}

Comparative and cross-cultural work complicates the picture in ways that are essential rather than inconvenient, because any account of harmonic information must be constrained at once by variation and by persistence. Studies of non-human animals suggest that some sensitivity to consonance-like or interval-structured organization may not be uniquely human, though the evidence is heterogeneous and should not be smoothed into a false consensus. Hoeschele and colleagues, Walker and colleagues, and related comparative work show that patterned sensitivity to harmonic organization can appear outside the human case, but also that such sensitivity varies by species, task, and ecological niche. These results are valuable precisely because they resist simplistic conclusions. They neither prove a universal biological law nor support the idea that harmonic organization is only a local artifact of Western music.

Cross-cultural work imposes an equally important discipline. The findings on the Tsimane have become especially important because they block any easy slide from recurring structure to universal aesthetic preference \parencite{mcdermott_2016}. The lesson is not that harmonic regularities disappear cross-culturally, but that familiarity, enculturation, and listening practice matter. This is one of the reasons HIT must keep distinguishing structural recurrence from evaluative judgment. A system may exhibit harmonic organization without every culture or listener assigning it the same aesthetic value. Cross-cultural evidence therefore does not invalidate the search for harmonic information; it forces the search to become more precise about what kind of universality is being claimed, and at what level.

Historical, ethnomusicological, and archaeological materials widen the frame further. Work on Aka and Mbuti polyphony, Tuvan throat singing, and other traditions grounded in overtone-rich or cyclic organization suggests that natural harmonic structures have been explored in multiple musical worlds, even when the local conceptual language differs from that of modern acoustics \parencite{arom_1991, turnbull_1961, levin_2006}. Archaeological and archaeoacoustic work, including discussions of prehistoric flutes and resonant spaces, likewise suggests that simple intervallic relations and resonant affordances are not late or marginal features of human culture \parencite{conard_2009, reznikoff_dauvois, morley}. None of this should be romanticized. The historical record does not authorize a fantasy of pristine universal harmony. It does show, however, that human engagement with structured resonant order is deep, diverse, and persistent.

What this comparative block contributes to HIT is therefore a double constraint. On one side, it keeps the framework from collapsing into a narrow Western theory of consonance. On the other, it keeps the critique of universalism from flattening every structured recurrence into pure contingency. Comparative and historical evidence belongs in the map because it shows that the problem of harmonic organization is neither culturally uniform nor culturally arbitrary. It is structured variation, and that is a more demanding object of inquiry than either universal harmony or radical relativism would allow.

\section{Information, prediction, and process organization}

The theory-of-information literature enters the map from a different angle, and without it patterned relation would remain descriptively interesting but conceptually underdetermined. It does not begin with consonance or resonance in the acoustic sense, yet it provides some of the clearest resources for understanding why patterned relations might matter informationally at all. Shannon's theory formalized information as reduction of uncertainty across a space of possible messages, giving modern science and engineering a rigorous way to measure communicative structure \parencite{shannon_1948}. That framework was deliberately abstract with respect to meaning, embodiment, and biological uptake. For HIT, this abstraction remains invaluable, but it is not sufficient by itself. It tells us how information can be measured, not yet why some relational patterns might have special organizational force.

Dubnov's work on information and musical expectation becomes important at exactly this point because it links structured predictability to informational salience in temporally organized patterns \parencite{dubnov_2004}. What matters in such cases is not randomness alone, but the relation between structure, surprise, and expectation. Friston's free-energy framework pushes this problem further into the domain of adaptive systems, suggesting that organisms remain viable in part by reducing uncertainty through better attunement to the structured regularities of their environment \parencite{friston_2010}. Even if one does not treat free energy as a universal master key, the convergence is evident: systems that track lawful structure effectively gain a predictive and organizational advantage over systems that do not.

Peirce adds a semeiotic dimension to this picture. If a sign is not simply an object but a triadic relation among sign, object, and interpretant, then information already appears as something irreducibly relational rather than merely additive or material \parencite{peirce_1931}. In this respect, the informational literature helps explain why HIT cannot be confined to bare acoustics or bare number. Ratios and recurrent couplings become relevant not because they are intrinsically sacred, but because they may support structured difference, prediction, and interpretation across systems. The strength of this literature lies in its ability to explain how patterned structure becomes informationally consequential. Its limit is that it often abstracts too far from the material regimes through which such structure is actually realized.

This is why the informational and predictive literatures belong in Chapter 6 rather than only in later theoretical chapters. They already participate in the convergence. They tell us that systems exploit structure, that prediction depends on lawful organization, and that information becomes meaningful through differential consequence. What they do not yet tell us, at least not on their own, is why simple harmonic relations so often reappear as the relevant kind of structure. That is the point at which they need the neighboring literatures on resonance, recurrence, and embodied coupling.

\section{Bioelectricity and biological organization beyond hearing}

One of the most important expansions of the map comes from biological organization outside the narrow auditory frame, because HIT weakens if it quietly collapses back into a theory of hearing. Levin's work on bioelectricity is exemplary here because it treats living pattern formation as sensitive to structured gradients, conductance ratios, and distributed electrical organization rather than only to local molecular identity \parencite{levin_2021}. The relevance of this work to HIT lies in its insistence that biological form can depend on relational electrical states and not merely on discrete material parts. This does not yet amount to a theory of harmonic information, but it places the question of ratio-like organization inside developmental and physiological regulation rather than restricting it to hearing or music.

Other literatures on coherence and resonance in biological or molecular systems widen the frame further, though with varying degrees of stability and controversy. Work on quantum coherence in photosynthetic systems, for example, suggests that living systems may exploit structured dynamical organization in ways that are not exhausted by classical local descriptions \parencite{scholes_2010}. More recent discussions of resonance-like behavior in microtubules and related molecular structures move in a similarly suggestive direction, though here caution is especially necessary because the field is heterogeneous and the interpretations are still in motion \parencite{celardo_2023}. These literatures should not be overread. They do not prove that life is organized by harmony in any simple sense. They do, however, broaden the empirical horizon in a way that matters for HIT: they show that questions of coherence, patterned coupling, and distributed relational organization already reach well beyond the sensory domain.

The main contribution of this literature is therefore strategic and substantive at once. It prevents the manuscript from quietly collapsing back into a theory about auditory consonance alone, and it shows that relational organization in living systems can become relevant at scales ranging from development to cellular dynamics. Its main limitation is that it does not yet provide a unified vocabulary compatible with the more formal accounts of recurrence and mode-locking discussed earlier. But that limitation is precisely why it belongs in a convergent map. HIT becomes stronger, not weaker, when it can show that the relevant clues already extend beyond the familiar auditory case.

\section{What each field knows, and what each field misses}

Table~\ref{tab:convergence-fields} summarizes the seven domains surveyed, the vocabulary each uses, and the kind of evidence each produces.

\begin{table}[ht]
\centering
\caption{Seven domains of empirical convergence.}
\label{tab:convergence-fields}
\small
\begin{tabular}{@{}lp{3.3cm}p{2.8cm}l@{}}
\toprule
Domain & Key finding & Local vocabulary & Mechanism \\
\midrule
Neuroscience & Neural resonance, CFC & Entrainment, coupling & Phase locking \\
Psychoacoustics & Consonance hierarchies & Roughness, harmonicity & Spectral interaction \\
Acoustic ecology & Niche partitioning & Soundscape & Spectral avoidance \\
Nonlinear dynamics & Mode locking at simple ratios & Synchronization & Arnold tongues \\
Comparative & Cross-species preferences & Perception & Template matching \\
Information theory & Prediction via structure & Compression, entropy & Redundancy reduction \\
Bioelectricity & Voltage pattern organization & Bioelectric code & Gap junctions \\
\bottomrule
\end{tabular}
\end{table}

The convergent picture can now be stated more clearly. Neuroscience shows that relations among frequencies and oscillatory bands are implicated in integration, entrainment, and cortical organization, but it often keeps those findings within the horizon of neural mechanism. Psychoacoustics and tonal perception show with unusual precision how consonance, roughness, harmonicity, and pitch depend on structured relations, yet they often treat the auditory case as the natural border of the problem. Acoustic ecology and biosemiotics show that patterned organization can be ecological, embodied, and semiotic, though they often lack the formal tools needed to describe recurrence and locking with mathematical sharpness. Nonlinear dynamics provides that formal sharpness in abundance, but usually without a commensurate account of perception, meaning, or living significance.

Comparative and historical evidence adds a crucial corrective by showing that harmonic organization is neither a simple universal taste nor a purely local fabrication. Cross-cultural findings restrain grandiosity; ethnomusicological and archaeological findings restrain relativistic flattening. The informational and predictive literatures clarify why structured relation matters at all, but they often abstract away from the concrete resonant media through which such structure becomes operative. Bioelectric and broader biological literatures extend the field of relevance beyond hearing, while still lacking a settled common language with the other domains. Each literature, then, sees a real part of the problem. None sees the whole.

This is the central conclusion of the chapter. The empirical convergence required by HIT is already present, but it is present in distributed form. The fields do not fail because they are weak. They fail only if one asks them, one by one, to articulate a common research object that was never their local mandate. Once they are read together, a more coherent picture emerges. Structured relations among oscillatory, resonant, recurrent, and mode-organized processes matter across physical, neural, ecological, perceptual, and biological domains. The meanings, mechanisms, and scales differ. The recurrence of the problem does not.

That recurrence is sufficient to justify the next step, but not to settle it. A convergent map is not yet a method. The fact that many literatures point toward a common object does not tell us how to compare descriptors, how to define evidence, how to distinguish natural-harmonic structure from perceptual or cultural transforms, or how to test claims causally and cross-modally. Those are methodological questions, and they require a more explicit framework than the literature itself can provide. That is the work of the next chapter.
